\section{Полугруппа}

\begin{definition}[Полугруппа]
    Множество $S$ с заданной на нём ассоциативной бинарной операцией $\cdot: S \times S \to S$ называется \emph{полугруппой} и обозначается $(S, \cdot)$.
    Если операция $\cdot$ является коммутативной, то говорят о \emph{коммутативной полугруппе}.
\end{definition}

\begin{example}
    Приведём несколько примеров полугрупп.
    \begin{itemize}
        \item Множество положительных целых чисел с операцией сложения является коммутативной полугруппой.
        \item Множество целых чисел с операцией взятия наибольшего из двух ($\max$) является коммутативной полугруппой.
        \item Множество всех строк конечной длины без пустой строки%
              \sidenote{
                  Множество всех строк конечной длины c пустой строкой также является полугруппой.
                  Однако, такая структура является ещё и моноидом, что будет показано далее.}
              над фиксированным алфавитом $\Sigma$ с операцией конкатенации является полугруппой.
              Так как конкатенация на строках не является коммутативной операцией, то и полугруппа не является коммутативной.
    \end{itemize}
\end{example}

